\documentclass[10pt,a4paper]{amsart}
\usepackage[margin=19mm]{geometry}
\usepackage{amsmath,amssymb,mathtools}
\usepackage[scr=rsfs,cal=euler]{mathalfa}
\usepackage{xfrac}
\usepackage[unicode,hidelinks,bookmarks=false]{hyperref}
\newcommand{\ie}{i.e.\ }
\newcommand{\cf}{cf.\ }
\newcommand{\resp}{respectively\ }
\newcommand{\Subsection}{\subsection}
\providecommand{\nobreakdash}{\nobreak\hbox{-}}
\setlength{\emergencystretch}{2em}
\multlinegap0pt
\usepackage{amssymb}
\usepackage[shortlabels]{enumitem}
\usepackage{mathtools}
\usepackage{xcolor}
\newtheorem{lemma}{Lemma}
\newtheorem{Cases}{Cases}
\newtheorem{definition}{Definition}
\newtheorem{claim}{Claim}
\newtheorem{example}{Example}
\newcommand\OT{\mathsf{OT}}
\newcommand{\Om}{{\Omega}}
\newcommand{\al}{{\alpha}}
\newcommand{\de}{\delta}
\newcommand{\ga}{\gamma}
\newcommand{\kp}{[k\leftarrow k+1]}
\newcommand{\la}{\lambda}
\newcommand{\ok}{o_{k}}
\newcommand{\okpe}{o_{k+1}}
\newcommand{\om}{\omega}
\def\scbr#1{\scbrl#1\scbrr}
\def\scbrl{[\hspace{-.14em}[}
\def\scbrr{]\hspace{-.14em}]}
\title{\texorpdfstring{Goodstein at the Second Threshold:\\ An independence Result for $\sf ID_2$}{Goodstein at the Second Threshold: An Independence Result for ID2}}
\author{Oriola Gjetaj, Andreas Weiermann}
\date{Source: arXiv:2604.00771v1 (CC BY 4.0)}
\begin{document}
\maketitle

\noindent\emph{Source and licence.} \texorpdfstring{Goodstein at the Second Threshold:\\ An independence Result for $\sf ID_2$}{Goodstein at the Second Threshold: An Independence Result for ID2}. Version arXiv:2604.00771v1 (CC BY 4.0), distributed by arXiv under the Creative Commons Attribution 4.0 International licence, http://creativecommons.org/licenses/by/4.0/, as recorded in the arXiv licence metadata for this exact version and shown on the abstract page https://arxiv.org/abs/2604.00771v1. Full licence notice: \texttt{licenses/arxiv-2604.00771-CC-BY-4.0.txt}.\par\smallskip

\noindent\textit{Editorial scope.} Counted unit: the article's lemma lm:LemmaInd (source0.tex lines 2291--2293). The article's complete original proof of it is delivered with it (source0.tex lines 2294--2605, one \texttt{proof} environment of 312 lines). No mathematics is written, repaired or re-derived: the statement and the proof are the article's own lines, in the article's own notation, and no retained line is substituted or re-flowed.  Retained uncounted as the source-local context the proof needs: lines 1437--1444; lines 1451--1459; lines 1482--1492; lines 2063--2069; lines 2103--2108.  Nothing was omitted from any retained span, so the package carries no omission record.  No retained line contains a citation, so the wrapper carries no bibliography and no bibliography entry.  No retained line contains a figure environment, an image-inclusion command or drawing code, so no image is delivered and the package declares no asset dependency.  Editorial formatting only, in addition: an AMS \texttt{amsart} preamble that restates the article's own declarations and macros used by the retained lines, namely the environments \texttt{lemma}, \texttt{Cases}, \texttt{definition}, \texttt{claim}, \texttt{example} on separate counters, together with the standard packages those lines need.  The retained environments print in this document as Lemma 1, Cases 1, Definition 1, Claim 1, Example 1 in order of appearance, whereas the article prints the same items with the numbering of its own full document.  The complete native source of the article as distributed in the arXiv e-print for this exact version is bundled unchanged as \texttt{sources/197611/source0.tex}.
\begin{lemma}\label{lm:HardyNF}
Fix $k\geq 2$.  For all $m\geq k$, there exist unique $\al\in  \OT_0\cap \Om$ and $q<\om$ such that
\begin{enumerate}
    \item $m=H_\al(k) +q$,
    \item $\al$ is maximal with $H_\al(k) \leq m$. %\revise{such that $H_\al(k) \leq m < H_{\al+1}(k)$.}\david{That last part is redundant since $\al$ is already maximal.}
    \item $q$ is the natural number with $q < H_{\al+1}(k)$ such that $H_\al(k) +q \leq m < H_\al(k) +q+1$.
\end{enumerate}
\end{lemma}

\begin{lemma}\label{lm:NForms} Fix $k\geq 2$.
\begin{enumerate}
    \item If $m=_{k} H_\al(k) \cdot p$ then $H_\al(k) \cdot (p-1)$ is in $k$ normal form as well.
    \item If  $m=_{k} H_\al(k) +q+1$ then $H_\al(k) + q$ is in $k$ normal form.
    \item If $m=_{k}H_\al(k)$ and $\al$ is a limit, then $H_{\al\{b\}}(k)$ for $0<b<k$ is in $k$ normal form. 
    \item If $m=_{k}H_{\al}(k)$ and $\al =\al_0+l$ is successor for $l$ a natural number,
    then $H_{\al_0+l-1}(k)$ is in $k$ normal form. 
\end{enumerate}
\end{lemma}

\begin{lemma}\label{lm:OrdMeasure}
Assume that  $\al$ is maximal with $H_\al(k)\leq m$. Then for all contexts $\la$,
\begin{enumerate}
    \item If $\al = \la\scbr{t}$ and $\al<\de= \la^-\scbr{\om}$ and $\al< \de[t+1]$  then $t <H_{\de\{k-1\}}(k)$.
    
    \item If $\al =\la\scbr{\ga}$ and $\al< \de=\la^-\scbr{\psi_0(\mu\scbr{\Om})}$ and if there is some $t$ such that
    $\de[t]\leq\al<\de[t+1]$ then $t<H_{\de\{k-1\}}(k)$.
    
    \item If $\al = \la\scbr{\psi_i(\ga)\cdot s}$ and $\al < \de = \la^-\scbr{\psi_i(\ga+1)}$  for $i \in \{0,1\}$ and $\al < \de[s+1]$ then $s< H_{\de\{k-1\}}(k)$. 
\end{enumerate}
\end{lemma}

\begin{lemma} \label{lm:okpe=ok} Let $m$ be in $k$ normal form and $m\kp$ the $k+1$ normal form after applying base change operation to $m$. Then the following hold;
\begin{enumerate}
    \item $\okpe(m') = \ok(m)$.
    \item $\okpe(\al') =\ok(\al)$.
    \item $\okpe((\la\scbr{\cdot})')=\ok(\la\scbr{\cdot})$.
\end{enumerate}
\end{lemma}

\begin{lemma} \label{lm:OkMon} Monotonicity  of ordinal assignment. %in system $\OT$.
\begin{enumerate}
    \item If $m<n$ then $\ok m < \ok n$. 
    \item If $\alpha<\beta$ then $\ok \alpha<\ok \beta$.
\end{enumerate}
\end{lemma}

\begin{lemma}\label{lm:LemmaInd} Let $m':= m\kp$. Given $k,m<\om$ with $k\geq 2$,
\[\okpe(m'-1) \geq \ok(m)[k].\]
\end{lemma}

\begin{proof}
We prove the claim by induction on $m$. 
The cases $m<k$ and $m=k$ are immediate from the definition of $o_k$. For $m=H_\alpha(k)\cdot p+q$ in $k$–normal form we split into $q>0$, $q=0<p$, and $q=0=p=1$, and in each case apply Lemmas \ref{lm:HardyNF}-\ref{lm:OrdMeasure} and \ref{lm:okpe=ok}-\ref{lm:OkMon} to track how base–change interacts with the Hardy hierarchy and with fundamental sequences.
\\
If  $m<k$ then $\okpe(m'-1) = \okpe(m-1)=m-1= \ok(m)[k]$.
\\
If $m=k$ then $m= H_0(k)$ and 
\begin{align*} 
    \okpe(m'-1)& =\okpe(H_{0} (k+1) -1)\\
    &= \okpe(k ) =k \\
    &= \psi_0(0)[k]=\ok(m)[k].
\end{align*}

Let $m=_k H_\al(k)\cdot p +q$.
\begin{Cases}
\item If $q>0$, then by Lemma \ref{lm:NForms}, $H_\al(k) \cdot p+ q-1$ is in $k$ normal form. Thus we can apply the ordinal assignment. By induction hypothesis $ \okpe(q'-1) \geq \ok(q)[k]$ and by Lemma \ref{lm:okpe=ok}, $\okpe(\al') = \ok(\al)$.
%In the third step, we apply the fundamental sequence to $\ok(q)$ since $\psi_0(\ok(\al))> \ok(q)$. % if this is not the case, then it means we have the number of ocurences of $\al$.
\begin{align*} 
        \okpe(m'-1)& =\okpe(H_{\al'} (k+1) \cdot p + q'-1)\\
        & =\psi_0(\okpe(\al'))\cdot p + \okpe(q'-1)\\
        &\geq \psi_0(\ok(\al))\cdot p + \ok(q)[k]=\ok(m)[k].
\end{align*}

\item If $q=0$ and $p>1$ then by induction hypothesis in step 3 to 4, as well as by Lemma \ref{lm:okpe=ok}, $\okpe(\al')=\ok(\al)$, we have the following,
    \begin{align*} 
            \okpe(m'-1)& =\okpe(H_{\al'}(k+1) \cdot p - 1)\\
            &= \okpe(H_{\al'}(k+1) \cdot (p-1) + H_{\al'}(k+1)-1)\\
            & =\psi_0(\okpe(\al')) \cdot (p-1) + \okpe(H_{\al'}(k+1)-1)\\
            &\geq_{ih}  \psi_0(\ok(\al))\cdot (p-1)  + \ok(H_\al (k))[k]\\
            &= \ok(H_\al(k)\cdot (p-1) + H_\al (k))[k] \\
            &=\ok(H_\al(k)\cdot p)[k].
        \end{align*}      
\item If $q=0$ and $p=1$ then $m=_k H_\al(k)$ where $\al = \psi_0\al_0 +... +\psi_0\al_n+l$, for simplicity we will write, $\al = \xi + l$.
There are several sub-cases to consider depending on the type of $\al$:
\begin{Cases}
        
        \item ($\al$ is a successor, $\al = \xi + l, 0<l<k$). Then $\ok(\al)$ is a successor; $\ok(\al)= \ok(\xi)+l $. Since $H_{\xi'+l-1}(k+1) \cdot k $ is in $k+1$ normal form by Lemma \ref{lm:NForms}, we can apply the ordinal assignment in steps 3 to 4.\\ Additionally by Lemma \ref{lm:okpe=ok}, $\okpe(\xi')=\ok(\xi)$.
        \begin{align*} 
            \okpe(m'-1)& =\okpe(H_{\xi'+l} (k+1) -1)\\
            & = \okpe(H_{\xi'+l-1} (k+1) \cdot (k+1) -1)\\
            & \geq \okpe(H_{\xi'+l-1}  (k+1)\cdot k )\\
            & = \psi_0({\okpe(\xi') +l-1}) \cdot k \\
            &= \psi_0({\ok(\xi)} + l-1) \cdot k\\
            &= \psi_0(\ok(\xi)+ l)[k] =\ok(m)[k]\\
        \end{align*}
        
        \item ($\al$ is a successor, $\al=\xi+l, l\geq k$). Then $l$ is written in $k$ normal form. By induction hypothesis, since $l<m$ then $\okpe(l'-1)\geq \ok(l)[k]$. Also, in $\OT$, $\psi_0(\okpe(l'-1)) \geq \psi_0(\ok(l)[k])$.\\
        Since the cofinality of $\ok(l) =\om$, we can apply the fundamental sequence to $\ok(l)$ in the last step. 
        By Lemma \ref{lm:okpe=ok}, $\okpe(\xi')=\ok(\xi)$.
        \begin{align*} 
            \okpe(m'-1)& =\okpe(H_{\xi'+l'} (k+1) -1)\\
            & \geq \okpe(H_{\xi'+l'-1}  (k+1) )\\
            & = \psi_0({\okpe(\xi')} + \okpe(l'-1)) \\
            &\geq \psi_0({\ok(\xi)}+ \ok(l)[k])\\
            &=\psi_0({\ok(\xi)}+ \ok(l))[k]=\ok(m)[k]\\
        \end{align*}
        
        
        \item ($\al$ is a limit of countable cofinality). Hence $\al = \xi + \psi_0(\al_n)$ where $\al_n$ can be $0$, a successor, of countable cofinality or of uncountable cofinality. We will differentiate on the type of $\al_n$ and $\ok(\al_n)$ in the following sub-cases:\medskip
        
        \begin{Cases}
            \item ($\al_n = 0$). Thus $\al = \xi + \om$ is of countable cofinality and $\ok(\al)=\ok(\xi)+\Om$ is of uncountable cofinality.
            Since $H_{\al\{d\}}(k)$ for $d<k$ is in $k$-normal form by Lemma \ref{lm:NForms} then we can apply the ordinal assignment.
            By Lemma \ref{lm:okpe=ok}, $\okpe(\xi')=\ok(\xi)$.
            \begin{align*} 
                \okpe(m'-1) & =\okpe(H_{\xi'+\om}(k+1) -1)\\
                & = \okpe(H_{(\xi'+\om)\{k+1\}}(k+1) -1)\\
                & \geq  \okpe(H_{(\xi'+\om)\{k\}}(k+1))\\
                %& \geq_{(Claim \ref{claim0})} \psi_0( \ok(\xi)+ \Om[z_{k}])\\
                & \geq_{(Claim \ref{claim0})} \psi_0( \ok(\xi)+ \Om )[k]\\
                &=\ok(H_{\xi+\om}(k)[k].\\
            \end{align*}
            We prove the following Claim \ref{claim0} to explain steps 3 to 4.
            First, we will define the fundamental sequences for our case.\\
            
            \begin{definition}Fundamental sequences for $\psi_0(\ok(\xi) + \Om)[t] = \psi_0(\ok(\xi) +\Om[z_t])$: \\
                $z_0 := 0$\\
                $z_1 := \psi_0( \ok(\xi)+ \Om)[0] = \psi_0( \ok(\xi)+ \Om[z_0]) = \psi_0( \ok(\xi) +z_0)=\psi_0(\ok(\xi)) $\\
                $z_{t+1} := \psi_0(\ok(\xi) + \Om)[t] = \psi_0(\ok(\xi) + \Om[z_t])=
                \psi_0( \ok(\xi)+ \Om[\psi_0(\ok(\xi) +  \Om[z_{t-1}])])
                = \psi_0( \ok(\xi)+ \psi_0(\ok(\xi) +  z_{t-1}))$.\\
            \end{definition}
            \begin{claim}\label{claim0}
                \[ \okpe(H_{(\xi'+\om) \{t\}}(k+1))\geq \psi_0( \ok(\xi)+ \Om)[t].\]
            \end{claim}
            \begin{proof} We prove the claim by induction on $t$.\\
                We use the fact that $ H_{(\xi'+\om)\{t\}}(k+1)$ is in $k+1$ normal form for $t<k+1$ by Lemma \ref{lm:NForms}.\\ %N\OTE
                For $t=0$, by Lemma \ref{lm:okpe=ok}, $\okpe(\xi')= \ok(\xi)$, and in $\OT$ we have that $ \psi_0(\okpe(\xi')) =  \psi_0(\ok(\xi))$.
                \begin{align*}
                    \okpe(H_{(\xi'+\om)\{0\}} (k+1)) &=\okpe(H_{(\xi'+\om)[0]}(k+1)) \\
                    &=\okpe(H_{\xi'+\om[0]}(k+1))\\
                    &= \okpe(H_{\xi'}(k+1))\\
                    &= \psi_0(\okpe(\xi')) = \psi_0( \ok(\xi)) \\
                    &= \psi_0( \ok(\xi)+ \Om[z_0]) = \psi_0( \ok(\xi)+ \Om)[0].\\
                \end{align*} 
                
                Since $t=0$ takes the argument to $0$, we will also consider at $t=1$.
                Again using  Lemma \ref{lm:okpe=ok}, $ \psi_0(\okpe(\xi') + \psi_0(\okpe(\xi'))) =  \psi_0(\ok(\xi) + \psi_0(\ok(\xi)))$.
                \begin{align*}
                    \okpe(H_{(\xi'+\om)\{1\}} (k+1)) &=\okpe(H_{\xi'+\om[H_{(\xi' + \om)\{0\}}(k+1)]}(k+1)) \\
                    &= \okpe(H_{\xi'+H_{\xi'}(k+1)}(k+1))\\
                    &= \psi_0(\okpe(\xi') + \psi_0(\okpe(\xi')))\\
                    &= \psi_0( \ok(\xi)+ \psi_0(\ok(\xi))) \\
                    &= \psi_0( \ok(\xi)+\Om[\psi_0(\ok(\xi) +\Om[z_0]))]) \\
                    &= \psi_0( \ok(\xi)+ \Om[z_1]) = \psi_0( \ok(\xi)+ \Om)[1].\\
                \end{align*}
                Induction step for $t+1$;
                \begin{align*}
                    \okpe(H_{(\xi'+\om)\{t+1\}} (k+1)) &=\okpe(H_{\xi'+\om[H_{(\xi'+\om)\{t\}}(k+1)]}(k+1)) \\
                    &= \okpe(H_{\xi'+H_{(\xi'+\om)\{t\}}(k+1)}(k+1)) \\
                    &=\psi_0(\okpe(\xi')+ \okpe(H_{(\xi'+\om)\{t\}}(k+1)))\\
                    &\geq \psi_0( \ok(\xi)+\psi_0(\ok(\xi)+ \Om)[t]) \\
                    &= \psi_0( \ok(\xi)+\Om[\psi_0(\ok(\xi)+ \Om[z_t])]) \\
                    &= \psi_0( \ok(\xi)+ \Om[z_{t+1}] ) = \psi_0( \ok(\xi)+ \Om)[t+1].\\
                \end{align*}
            \end{proof}
            
            { \em The following cases, $\al=\xi+\psi_0(\ga+l)$, $\al = \xi+\psi_0(\de+\psi_0(\ga+l))$, $\al= \xi+\psi_0(\de + \psi_1(\ga+l))$ and so on, have the same procedure, hence we will consider them together in the following case, and use a context to write $\al$.} \\
            
            \item \label{CaseSucc}($\al_n = \ga+l$ or $\la^-\scbr{\psi_i(\ga+l)}$, for $i\in \{0,1\}$, and $0<l<k$ and  $\ga\in Lim$ or $0$).\smallskip
            \\
            Then $\al= \xi+ \la^-\scbr{\psi_i(\ga+l)}$ and $\ok(\al) = \ok(\xi) + \ok(\la^-)\scbr{\psi_{i+1}(\ok(\ga)+l)}$ are both of countable cofinality.
            \\
            \begin{example}
                For example; \\
                For $\al=\xi+\psi_0(\ga+l)$, the context is 
                $\la^-\scbr{\cdot}=\scbr{\cdot}$ and $\al=\xi +\la^-\scbr{\psi_0(\ga+l)}$.
                \\
                For $\al = \xi+\psi_0(\de+\psi_0(\ga+l))$, the context is $\la^-\scbr{\cdot} = \psi_0(\de + \scbr{\cdot})$, then $\al= \xi +\la^-\scbr{\psi_0(\ga+l)}$.
                \\
                If $\al = \xi + \psi_0(\de + \psi_1(\ga+l))$, the context is $\la^-\scbr{\cdot} = \psi_0(\de + \scbr{\cdot})$, then
                $\al= \xi+ \la^-\scbr{\psi_1(\ga+l)}$.
                \medskip
                \\
                The following example illustrates the ordinal assigment:\\
                $\ok(H_{\psi_0\psi_0\psi_01})[k]= \psi_0(\psi_1\psi_1\psi_11)[k] = \psi_0(\Om^3\cdot \om)[k]= \psi_0(\Om^3 \cdot k)$.
                \\
                Consider another example to ordinal assignment as in the following: \\
                $\ok(H_{\psi_0(\psi_1\Om + \psi_1(\om+1))})[k] = \psi_0(\psi_1(\psi_2\Om_2+ \psi_2(\Om+1)))[k]\\=
                \psi_0(\psi_1(\psi_2\Om_2 + \psi_2\Om\cdot k))$. 
            \end{example}
            
            By Lemma \ref{lm:NForms} we have that
            $H_{\xi'+(\la^-)'\scbr{\psi_i(\xi'+l-1)\cdot k}}(k+1)$ is in $k+1$ normal form, hence in steps 3 to 4 we can apply the ordinal assignment.
            \begin{align*} 
                \okpe(m'-1) & =\okpe(H_{\xi'+(\la^-)'\scbr{\psi_i(\ga'+l)}}(k+1) -1)\\
                & = \okpe(H_{(\xi'+(\la^-)'\scbr{\psi_i(\ga'+l)})\{k+1\}}(k+1) -1)\\
                & \geq  \okpe(H_{\xi'+(\la^-)'\scbr{\psi_i(\ga'+l-1)\cdot k}}(k+1))\\
                & =  \psi_0(\okpe(\xi')+ \okpe((\la^-)')\scbr{\psi_{i+1}(\okpe(\ga') + l-1)\cdot k})\\
                & = \psi_0( \ok(\xi)+ \ok(\la^-)\scbr{\psi_{i+1}(\ok(\ga) +l-1) \cdot k})\\
                %& = \psi_0( \ok(\xi)+ \ok(\la^-)\scbr{\psi_{i+1}(\ok(\ga) +l)}[k] )\\
                & = \psi_0( \ok(\xi)+ \ok(\la^-)\scbr{\psi_{i+1}(\ok(\ga) +l)} )[k]\\
                &=\ok(H_{\xi+\la^-\scbr{\psi_i(\ga+l)}}(k))[k].\\
            \end{align*}		
            
            
            \item \label{CaseNotSucc} $(\al_n=\ga+l$ or $\mu^-\scbr{\psi_i(\ga+l)}$ for $i \in \{0,1\}$,  $l\geq k$ and $\ga\in Lim$ or $0$).\\
            As above we will cover several cases with the same procedure and we will use a context $\mu^-\scbr{\cdot}$ to write $\al$. Since $l\geq k$, then also $l$ is written in $k$  normal form. Then $\al = \xi + \mu^-\scbr{\psi_i(\ga+l)}$ and $\ok(\al)= \ok(\xi)+ \ok(\mu^-)\scbr{\psi_{i+1}(\ok(\ga)+\ok(l))}$.
            \\
            \begin{example}
                Here we cover ordinals such as $\al = \xi + \psi_0(\ga+l)$ or $\al =\xi+ \psi_0(\de+\psi_0(\ga+l))$ or $\al = \xi+ \psi_0(\de + \psi_1(\ga+l))$. 
                %The reason we write the context differently here compared to the case above is that the inner $\psi$ term was important whereas here the differences are erased.
                \\
                \\
                As for the ordinal assignment, consider an example for $k=3$ and $\al =\psi_0\psi_0\psi_03$: \\
                $\ok(H_{\psi_0\psi_0\psi_03})[k]= \psi_0(\psi_1\psi_1\psi_1\om)[k] = \psi_0(\Om^3\cdot \om^\om)[k]= \psi_0(\Om^3\cdot \om^{k}).$
                \\
                Or consider $\al = \psi_0(\psi_1(\om + 3 ))$, such that\\ 
                $\ok(H_{\psi_0(\psi_1(\om + 3 ))})[k] = \psi_0(\psi_1( \psi_2(\Om + \om)))[k]  = \psi_0( \psi_1(\psi_2(\Om + \om[k]))).$
                \\
            \end{example}
            In step 3, by Lemma \ref{lm:NForms} we can apply the ordinal assignment, 
            \\as $ H_{\xi'+(\mu^-)'\scbr{\psi_i(\ga'+l'-1)}}(k+1)$ is in $k+1$ normal form. \\
            In steps 4 to 5, by induction hypothesis $\okpe(l-1)\geq \ok(l)[k]$. Additionally
            $\psi_{i+1}(\okpe(l-1)) \geq \psi_{i+1}(\okpe(l)[k])$ in $\OT$.\\
            Since the cofinality of $\ok(l)=\om$ we can apply the fundamental sequence to $\ok(l)$.
            
            \begin{align*} 
                \okpe(m'-1) & =\okpe(H_{\xi'+(\mu^-)'\scbr{\psi_i(\ga'+l')}}(k+1) - 1)\\
                & = \okpe(H_{\xi'+(\mu^-)'\scbr{\psi_i(\ga'+l')}\{k+1\}}(k+1)-1)\\
                %&= \okpe(H_{\xi'+\la'\scbr{\psi_0(\ga'+l')\{k+1\}}}(k+1) -1)\\
                & \geq \okpe(H_{\xi'+(\mu^-)'\scbr{\psi_i(\ga'+l'-1)}}(k+1))\\
                & = \psi_0(\okpe(\xi')+ \okpe((\mu^-)')\scbr{\psi_{i+1}(\okpe(\ga')+\okpe(l'-1))}))\\
                & \geq \psi_0( \ok(\xi)+ \ok(\mu^-)\scbr{\psi_{i+1}(\ok(\ga)+\ok(l)[k])})\\
                & = \psi_0(\ok(\xi)+ \ok(\mu^-)\scbr{\psi_{i+1}(\ok(\ga)+\ok(l))})[k]\\
                &=\ok(H_{\xi+ (\mu^-)\scbr{\psi_i(\ga+l)}}(k))[k].\\
            \end{align*}
            
            
            
            \item ($\al_n=\tau^-\scbr{\om}$).\label{CaseCountCof}
            Then $\al= \xi + \psi_0(\tau^-\scbr{\om})$ is of countable cofinality and
            $\ok(\al)= \ok(\xi) + \psi_1(\ok(\tau^-)\scbr{\Om})$ is of uncountable cofinality.\medskip
            \\
            Here we  cover cases such as  $\al = \xi + \psi_0(\de + \om)$ or $\al = \xi + \psi_0(\de+ \psi_0(\ga +\om))$  or $\al = \xi + \psi_0(\de+ \psi_1(\ga +\om))$. \\	    
            Recall the fundamental sequences for $\al = \psi_0(\al_0)$ where $tp(\al)=\om$ and $\al_0 :=\mu^-\scbr{\Om}$ for a $\psi_0-$free context $\mu^-$.
            $\psi_0(\al_0)[x] = \psi_0(\al_0[z_x])$ where $z_0=0$ and $z_{x+1}= \psi_0(\al_0[z_x])$.\medskip
            
            \begin{example}
                The following examples clarify the fundamental sequences that arise in this setting:\\
                For simplicity, let $\al = \psi_0(\om)$, then $\ok(H_{\psi_0(\om)}(k))[x]= \psi_0(\psi_1(\Om))[x]= \psi_0(\psi_1(\Om)[z_x])$ since $\psi_1(\Om)$ is of cofinality $\Om$. Then $\psi_0(\psi_1(\Om)[z_x])= \psi_0(\psi_1(\Om[z_x]))$. \\
                The following example illustrates further:\\
                Let $\al = \psi_0(\psi_1(\om))$ then $\ok(H_{\psi_0(\psi_1(\om))}(k))[x] = \psi_0(\psi_1(\om^{\Om_2+\Om}))[x]= \psi_0(\psi_1(\om^{\Om_2+\Om})[z_x])$ since cofinality of $\psi_1(\om^{\Om_2+\Om})$ is $\Om$.\\
                Then
                $ \psi_0(\psi_1(\om^{\Om_2+\Om})[z_x])= \psi_0(\psi_1(\om^{\Om_2+\Om}[z_x]))= \psi_0(\psi_1(\om^{\Om_2+\Om[z_x]}))$ since the cofinality of $\om^{\Om_2+\Om}$ is $\Om$.\\
            \end{example}
            
            Since $H_{\al\{d\}}(k)$ for $d<k$ is in $k+1$-normal form by Lemma \ref{lm:NForms} then we can apply the ordinal assignment.
            By Lemma \ref{lm:okpe=ok}, $\okpe(\xi')=\ok(\xi)$.
            
            \begin{align*} 
                \okpe(m'-1) & =\okpe(H_{\xi'+\psi_0((\tau^-)'\scbr{\om})}(k+1) -1)\\
                & = \okpe(H_{(\xi'+\psi_0((\tau^-)'\scbr{\om}))\{k+1\}}(k+1) -1)\\
                %&= \okpe(H_{(\xi'+\psi_0((\tau^-)'\scbr{\om}))\{k+1\}}(k+1) -1)\\
                & \geq  \okpe(H_{(\xi'+\psi_0((\tau^-)'\scbr{\om}))\{k\}}(k+1))\\
                %& \geq_{\text{Claim \ref{ClaimIndOm}}}\psi_0( \ok(\xi)+ \psi_1(\ok(\tau)\scbr{z_k }) )\\
                %&= \psi_0( \ok(\xi)+ \psi_1(\ok(\tau^-)\scbr{\Om}) [z_{k}])\\
                & \geq_{\text{Claim \ref{ClaimIndOm}}} \psi_0( \ok(\xi)+ \psi_1(\ok(\tau^-)\scbr{\Om})) [k]\\		
                &=\ok(H_{\xi+\psi_0(\tau^-\scbr{\om})}(k)[k].\\
            \end{align*}
            We prove Claim \ref{ClaimIndOm} to explain steps 4 to 5.
            First, we give the following definition.
            \begin{definition} Definition of fundamental sequences for 
                $\psi_0(\ok(\xi)+ \psi_1(\ok(\tau^-)\scbr{\Om}))[t]$;\\
                $z_0 =0$,\\
                $z_1 =	\psi_0( \ok(\xi)+ \psi_1(\ok(\tau^-)\scbr{\Om}))[0]                =\psi_0( \ok(\xi)+ \psi_1(\ok(\tau^-)\scbr{\Om})[z_0])
                =\psi_0( \ok(\xi) + \psi_1(\ok(\tau^-\scbr{\Om[z_0]}))$, 
                \\
                $z_{t+1} = \psi_0( \ok(\xi)+ \psi_1(\ok(\tau^-)\scbr{\Om}))[t]
                = \psi_0( \ok(\xi)+ \psi_1(\ok(\tau^-)\scbr{\Om})[z_t])\\
                =\psi_0( \ok(\xi)+ \psi_1(\ok(\tau^-)\scbr{\Om[z_t]}))
                = \psi_0( \ok(\xi)+\psi_1(\ok(\tau^-) \scbr{ \psi_0( \ok(\xi)+ \psi_1(\ok(\tau^-)\scbr{\Om[z_{t-1}]}) )}). $\medskip
            \end{definition}
            
            \begin{claim}\label{ClaimIndOm}
                \[\okpe(H_{(\xi'+\psi_0((\tau^-)'\scbr{\om}))\{t\}}(k+1))  \geq \psi_0( \ok(\xi)+ \psi_1(\ok(\tau^-)\scbr{\Om}))[t].\]
            \end{claim}
            \begin{proof}
                By induction on $t$. $H_{\xi'+\psi_0((\tau^-)'\scbr{\om})[x]}(k+1)$ is in $k+1$ normal form for $x<k+1$, hence we can apply the ordinal assignment. For $t=0$ we obtain
                \begin{align*}
                    \okpe(H_{(\xi'+\psi_0((\tau^-)'\scbr{\om}))\{0\}}(k+1)) &= \okpe(H_{\xi'+\psi_0((\tau^-)'\scbr{\om})[0]}(k+1)) \\
                    &= \okpe(H_{\xi'+\psi_0((\tau^-)'\scbr{0})}(k+1)) \\
                    &= \psi_0(\okpe(\xi')+\psi_1(\okpe((\tau^-)')\scbr{0}))\\
                    &= \psi_0( \ok(\xi)+ \psi_1(\ok(\tau^-)\scbr{0 }))\\
                    &= \psi_0( \ok(\xi)+ \psi_1(\ok(\tau^-)\scbr{\Om}))[0].\\
                \end{align*}    
                For $t=1,$
                \begin{align*}
                    \okpe(H_{(\xi'+\psi_0((\tau^-)'\scbr{\om}))\{1\}}(k+1)) 
                    &= \okpe(H_{\xi'+\psi_0((\tau^-)'\scbr{\om}) [H_{(\xi' + \psi_0((\tau^-)'\scbr{\om}))\{0\}}(k+1)])}(k+1)) \\
                    &=\okpe(H_{\xi'+\psi_0((\tau^-)'\scbr{H_{\xi' + \psi_0((\tau^-)'\scbr{0})}(k+1)})}(k+1)) \\
                    &= \psi_0(\okpe(\xi')+\psi_1(\okpe((\tau^-)')\scbr{\psi_0(\okpe(\xi')+\psi_1(\okpe((\tau^-)'\scbr{0}))}) )\\
                    &= \psi_0( \ok(\xi)+ \psi_1(\ok(\tau^-)\scbr{\psi_0(\ok(\xi)+ \psi_1(\ok(\tau^-\scbr{0}))) }))\\
                    &= \psi_0( \ok(\xi)+ \psi_1(\ok(\tau^-)\scbr{z_1}))\\
                    &= \psi_0( \ok(\xi)+ \psi_1(\ok(\tau^-)\scbr{\Om}))[1].
                \end{align*}
                Induction step for $t+1$ by the induction hypothesis;
                \begin{align*}
                    \okpe(H_{(\xi'+\psi_0((\tau^-)'\scbr{\om}))\{t+1\}}(k+1)) &=\okpe(H_{\xi'+\psi_0((\tau^-)'\scbr{\om})	[H_{(\xi' + \psi_0(\tau'\scbr{\om}))\{t\}}(k+1)]}(k+1)) \\
                    &= \okpe(H_{\xi'+\psi_0((\tau^-)'\scbr{ H_{(\xi'+\psi_0((\tau^-)'\scbr{\om}))\{t\}}(k+1)})}(k+1)) \\
                    &=\psi_0(\okpe(\xi')+ \psi_1(\okpe((\tau^-)') \\
                    &	\scbr{	\okpe( H_{(\xi'+\psi_0(\tau'\scbr{\om}))\{t\}})}))\\
                    &\geq \psi_0( ok(\xi)+\psi_1(\ok(\tau^-) \scbr{ \psi_0( \ok(\xi)+ \psi_1(\ok(\tau)\scbr{z_t}) )})) \\
                    &=\psi_0( \ok(\xi)+ \psi_1(\ok(\tau^-)\scbr{z_{t+1}}))\\
                    &=\psi_0( \ok(\xi)+ \psi_1(\ok(\tau^-)\scbr{\Om}))[t+1].
                \end{align*}
            \end{proof}
            

            
            \item ($\al_n = \mu^-\scbr{\Om}$) is of uncountable cofinality of the form $\psi_1(\ga_0)$ and $\ga_0$ is either 0 or of type $\Om$. The case where $\ga_0$ is a successor is covered in  \ref{CaseSucc} and \ref{CaseNotSucc}, and the case where $\ga_0$  is of countable cofinality is covered above in \ref{CaseCountCof}.
            We use the context to write $\al= \xi +\psi_0(\mu^-\scbr{\Om})$  of countable cofinality, then  $\ok(\al)= \ok(\xi) + \psi_1(\ok(\mu^-)\scbr{\Om_2})$ is of countable cofinality.
            We cover cases such as $\al = \xi + \psi_0(\Om)$ or $\al = \xi +\psi_0(\psi_1(\Om))$.\smallskip
            \\
            Recall the fundamental sequences for $\al = \psi_0(\al_0)$ where $tp(\al)=\om$ and $\al_0 :=\mu^-\scbr{\Om}$ for a $\psi_0-$free context $\mu^-$.
            $\psi_0(\al_0)[x] = \psi_0(\al_0[z_x])$ where $z_0=0$ and $z_{x+1}= \psi_0(\al_0[z_x])$.\medskip
            
            \begin{example}
            We now present an example to clarify the fundamental sequences involved in this case:\\
                Let $\al = \psi_0(\Om)$ then $\ok(H_\al(k))[x]= \psi_0(\psi_1(\Om_2))[x]= \psi_0(\psi_1(\Om_2)[x])$
                since $tp(\psi_1(\Om_2))=\om$.
                Instead, $\psi_0(\Om)[x]=\psi_0(\Om[z_x])$, where $z_0=0$ and $z_{x+1}=\psi_0(\psi_1(\Om_2)[z_x])$.
                This can be written fully as $\psi_0(\Om)[x] = \underbrace{\psi_0(\psi_0(\ldots\psi_00\ldots))}_{k \text{ times}}$.
                %= \psi_0(\psi_1(z_x))$.\\
                Whereas, $H_{\al\{x\}}(k)= H_{\psi_0(\Om)\{x\}}(k)\geq H_{\psi_0(\Om)[x]}(k).$\smallskip\\
                %= H_{\psi_0(\Om	[z_{H_{\psi_0(\Om)\{x-1\}}(k)}])}(k)$.
                The case described below represents a generalized instance of the preceding analysis.\medskip
            \end{example}
            We use the fact that $H_{\al'\{d\}}(k+1)$ is $k+1$ in normal form for $d<k+1$.
            As well as Lemma \ref{lm:okpe=ok} for $\okpe(\xi') = \ok(\xi)$.
            
            \begin{align*} 
                \okpe(m'-1) &= \okpe(H_{\xi'+\psi_0((\mu^-)'\scbr{\Om})}(k+1) -1)\\
                & = \okpe(H_{(\xi'+\psi_0((\mu^-)'\scbr{\Om}))\{k+1\}}(k+1) -1)\\
                & \geq  \okpe(H_{(\xi'+\psi_0((\mu^-)'\scbr{\Om}))\{k\}}(k+1))\\
                & \geq  \okpe(H_{(\xi'+\psi_0((\mu^-)'\scbr{\Om}))[k]}(k+1))\\
                &= \okpe(H_{(\xi'+\underbrace{\psi_0(
                (\mu^-)'[\psi_0((\mu^-)'[\ldots [\psi_0( (\mu^-)'  \scbr{0})] \ldots ])]
                ))}_{k \text{ times}}
                }(k+1))\\
                &= \psi_0(\okpe(\xi') + \underbrace{
                \psi_1(\okpe(\mu^-)'[\psi_1(\okpe(\mu^-)'[\ldots [\psi_1( \okpe(\mu^-)'  \scbr{0})] \ldots ])])}_{k \text{ times}})\\
                &= \psi_0(\ok(\xi) + \underbrace{
                \psi_1(\ok(\mu^-)[\psi_1(\ok(\mu^-)[\ldots [\psi_1( \ok(\mu^-) \scbr{0})] \ldots ])])}_{k \text{ times}})\\
                &= \psi_0( \ok(\xi)+ \psi_1(\ok(\mu^-)\scbr{\Om_2})[{k}] )\\
                &=\ok(H_{\xi+\psi_0(\mu^-\scbr{\Om})}(k))[k]\\
                &= \ok(m)[k].
            \end{align*}			
        \end{Cases}
    \end{Cases}
\end{Cases}
\end{proof}		

\end{document}
